\chapter{The Hard Problem Dissolved: A Complete Account}
\label{ch:hard_problem_dissolved}

\begin{chapterabstract}
This chapter revisits the hard problem of consciousness (Chapter~\ref{ch:hard_problem}) with the full theoretical apparatus now in hand. The problem dissolves not because we stop caring about subjective experience but because we have identified the physical process that subjective experience \emph{is}: categorical completion of oscillatory circuits in the \Hplus field. Each step of the dissolution is stated explicitly and mapped to its proof location in the monograph. The privacy axiom (Chapter~\ref{ch:geometry}) is shown to be a consequence of the non-grounded naming circuit structure (Chapter~\ref{ch:identity}), not an independent assumption. Qualia---the specific qualitative character of experiences---are selection signatures from the $\sim 10^6$ available weak-force completions, explaining why experiences have the particular character they do.
\end{chapterabstract}

The hard problem of consciousness, as Chalmers formulated it, comprises three distinct sub-problems:

\begin{enumerate}
  \item \textbf{The existence problem}: Why is there subjective experience at all, rather than mere information processing?
  \item \textbf{The character problem}: Why does experience have the particular qualitative character it has (the redness of red, the painfulness of pain)?
  \item \textbf{The privacy problem}: Why is conscious content accessible only to the subject and not directly to third-person observers?
\end{enumerate}

The categorical completion framework addresses all three.

\section{The Existence Problem}

Conscious experience exists because cascade termination causes cell death (Chapter~\ref{ch:oscillatory_holes}, Theorem~\ref{thm:cascade_termination}). Oscillatory holes must be filled. Neural circuit completion is not optional in oxygen-metabolising organisms operating above a threshold complexity; it is a thermodynamic necessity. The question ``why is there experience rather than mere processing?'' conflates two descriptions of the same process: experience \emph{is} processing, specifically the categorical completion process.

\section{The Character Problem}

The qualitative character of an experience is its selection signature: the specific one-of-$10^6$ weak-force completion chosen from the available configurations at that moment (Chapter~\ref{ch:oscillatory_holes}). Two experiences of ``the same'' object differ because the cytoplasmic state, neural connectivity, and historical BMD configuration differ, selecting different completions. This makes the character problem tractable: qualia are not primitive non-physical properties but physically realised selection events in partition space.

\section{The Privacy Problem}

Privacy is a structural consequence of the non-grounded naming circuit topology (Chapter~\ref{ch:identity}). A circuit without external ground cannot be probed from outside without becoming part of the circuit and thereby transforming what is being probed. This is formalised in the Privacy Axiom (Chapter~\ref{ch:geometry}) and is not a contingent feature of current measurement technology but a topological invariant of the conscious system's architecture.
